\chapter{Introduction}
\label{ch:Introduction}
\section{Background}

When the British archaeologist Arthur Evans in the beginning of the 20th century went to excavate Knossos on Crete, he found a huge amount of clay tablets written in a previously unknown linear script.
Some documents were clearly of a newer script, which he called \textit{Linear B}, while others were of an older one, which he named \textit{Linear A}.
Some fifty years later, the British amateur archaeologist Michael Ventris finally succeeded in deciphering the newer script Linear B, which he proved to be representing Mycenaean Greek, the earliest attested form of Greek.
But the same hypothesis did not hold for the older script Linear A, whose encoded language remained unknown, still yet to be deciphered.

A key obstacle to its decipherment has been the small size of its corpus, currently consisting of only around 1,500 documents - significantly smaller than that of its younger counterpart, which now spans more than 6,000 documents.\footcite{SteeleAegeanWriting2023} 
But new inscriptions are still being discovered, with the corpus most recently extended in 2025 by over 100 new documents.\footcite{CNRLinearACorpus2025}
Another challenge has been the lack of digital representations of these inscriptions, in a form suitable for statistical analysis.
While photographs with accompanying drawings of nearly the entire corpus are available online through \acf{Cefael} (see \textcite{EtudesCretoises21}), these are only available in print form.

This motivated the development of the palæographical\footnote{The term "palæographical" refers to the study of ancient writing systems.} database \textit{The Signs of Linear A} (SigLA), which seeks to "[develop] a systematic, exhaustive and user-friendly open access database of all Linear A inscriptions [...] in order to carry out statistical and palæographic analyses within the epigraphic corpus" \parencite{SalgarellaCastellanSigLA2020}.
However, the last addition to SigLA was in 2021, and the project has so far managed to document just 772 of all currently known documents.\footcite{SigLAAboutPage}\footcite{SigLASearchDocuments}

\subsection{A problem of manual digitization}

The drawings currently stored on SigLA were made by Ester Salgarella using the painting tool Krita.
For each document, she manually created her own annotated drawings based on the originals from \acs{Cefael}.
From these, she proceeded to create layered Krita files with metadata for each sign, which could then at last be imported to SigLA.\footcite{SalgarellaCastellanSigLA2020}
According to her (personal communication), this has been a highly time-consuming undertaking. 
And the necessity of this kind of labor-intensive work has arguably been the biggest reason to why SigLA still does not contain the entire available corpus.

The aim of this project is therefore to investigate whether parts of this process can be automated. More specifically, whether some kind of semi-automatic segmentation tool 
can accelerate the workflow, while maintaining high-quality outputs that are still compatible with SigLA's database.

\subsection{A problem of small data}
The small size of the Linear A corpus and the irregularities in quality of its documents is a huge problem for the application of deep learning approaches, which typically require large amounts of data to perform well.
Consequently, this project will instead seek to make use of generalized machine learning and computer vision, in particular Meta's state-of-the-art interactive segmentation tool, the \acf{SAM} \parencite{sam}, which is pre-trained on a massive and diverse dataset of over 1 billion masks.
\acs{SAM}'s generalization and interactivity (allowing user-placed markers that guide the model) might thus allow for semi-automatic annotation of Linear A inscriptions without the need for large amounts of training data, which is unfortunately not available in this context.\footnote{By the "interactivity" of \acs{SAM}, what is meant is NOT any built-in user interface in the form of a GUI, but rather the possibility of prompt-based placement of rectangles of focus, and points of additions and deletions that guide the model's prediction.}

\cleardoublepage
\section{Research questions}

\section{Outline of the thesis}